\chapter{Conclusions}
\label{chap:conclusions}

\begingroup
\setlength{\emergencystretch}{2em}

\noindent This chapter closes the traceability chain that begins with the healthcare
interoperability problem in Chapter~\ref{chap:introduction}, acquires its
analytical foundations in Chapter~\ref{chap:stateofart}, and produces empirical
evidence in Chapter~\ref{chap:selected-approach}. Its objective is to synthesise
the research contributions, reconcile the observations with the stochastic
invariants that delimit their interpretation, assess the limitations of the
available evidence, and derive a concrete engineering roadmap from the measured
behaviour. The resulting synthesis treats performance as a joint property of
latency distribution, useful completion, admission control, and workflow state:
no single scalar replaces this joint view, and no client-side measurement is
presented as an unsupported decomposition of a protected regional platform.

\section{Connecting Objectives to Evidence}
\label{sec:conclusions-objectives-evidence}
\subsection{Overview}
\noindent Chapter~\ref{chap:introduction} defines four objectives: interpret the
selected standards and architectures, formalise the corresponding healthcare
workflows, translate them into executable and observable workloads, and compare
direct with gateway-mediated paths. The Request Dataset Generator (RDG) provides
the artefact that connects these objectives. Its evidence chain preserves
configuration, seed, timeline, occurrence identity, request and flow outcomes,
and statistical population from workload construction to reporting. The three
results below provide the principal closure: path-level displacement
characterises the externally visible difference between the declared path
populations, stochastic truncation explains the different completion behaviour
of interrupted flows, and queueing saturation connects client-side delivery
drift with the deformation of useful completion under extreme offered load.
Table~\ref{tab:structural-reconciliation} provides the corresponding
objective--model--evidence map.

\colorlet{veryPeriLight}{veryPeri!18}
\begin{table}[H]
    \centering
    \caption[Structural Reconciliation Matrix]{Structural reconciliation
    matrix mapping the thesis objectives to the formal invariants of
    Chapter~\ref{chap:stateofart} and the empirical findings of
    Chapter~\ref{chap:selected-approach}. Source: Author's elaboration.}
    \resizebox{\textwidth}{!}{%
        \begin{tabular}{|
            >{\raggedright\arraybackslash}p{3.5cm}|
            >{\raggedright\arraybackslash}p{4.5cm}|
            >{\raggedright\arraybackslash}p{4.5cm}|
            >{\raggedright\arraybackslash}p{3.2cm}|}
            \hline
            \rowcolor{veryPeriLight}
            \textbf{Research Objective} &
            \textbf{Formal Invariant or Model} &
            \textbf{Empirical Evidence} &
            \textbf{System Status or Outcome} \\
            \hline
            \rowcolor{tableGray}
            Interpret the selected standards and architectures. &
            Occurrence-local dependency model and surrogate observability
            boundary. &
            Executable regional flow families preserve security, identity,
            document, and signature prerequisites across 74 complete live
            runs. &
            Reconciled within the Veneto integration perimeter; external
            validity remains scoped. \\
            \hline
            Formalise healthcare workflows. &
            Sequential intra-flow states and stochastic truncation
            $\mathcal{T}_{\tau}$ at the first unsuccessful transition. &
            The failed mediated-minus-direct population has a median
            differential of $+12.88$~ms; the invalid suffix receives no
            dispatch. &
            Fail-fast behaviour is supported; failed-flow tail dispersion
            remains visible. \\
            \hline
            \rowcolor{tableGray}
            Translate the workflows into executable and observable workloads. &
            Evidence-chain identity, planned-versus-actual dispatch, and a
            configured flow-concurrency bound. &
            The final dataset contains 757{,}517 structured events and
            158{,}449 flow completions; active flows plateau at 30 while drift
            grows under saturation. &
            Instrument validity holds below the admission boundary; overload
            exposes client-side backpressure. \\
            \hline
            Compare direct and gateway-mediated paths. &
            Path-level Overhead
            $\Delta L(P_{50})=L_{\mathrm{mediated}}-L_{\mathrm{direct}}$. &
            Successful-flow medians of $755.8$~ms and $571.1$~ms        
            $\Delta L(P_{50})=184.7$~ms. &
            The client-boundary displacement is quantified; mediation-only
            and component-level attribution remain non-assessable. \\
            \hline
        \end{tabular}%
    }
    \label{tab:structural-reconciliation}
\end{table}
\newpage
\subsection{Mediation Overhead and Path-Level Displacement}
\label{subsec:conclusions-path-displacement}

\noindent The fourth objective of Chapter~\ref{chap:introduction} requires a
differential measurement because internal gateway telemetry lies outside the
authorised observation perimeter. Chapter~\ref{chap:stateofart} consequently
defines the Path-level Overhead for a declared statistic $q$ as
\begin{equation}
    \Delta L(q)
    = L_{\mathrm{mediated}}(q)-L_{\mathrm{direct}}(q),
\end{equation}
subject to compatibility of transaction semantics, payload class, terminal
outcome, connection policy, offered load, and observation window. This
definition fixes the estimand before interpreting its magnitude.

\noindent The distinction between network, transport, and application-level completion
time developed by Iorio et al. determines the observation boundary. The client
observes the complete path effect, whereas the available evidence does not
isolate one internal processing stage\footcite{iorio-when-latency-matters}.

\noindent Chapter~\ref{chap:selected-approach} instantiates this estimand at the median.
For successful flows, the declared mediated population has a median completion
time of $755.8$~ms and the declared direct comparison population has a median
of $571.1$~ms. Their empirical displacement is
\begin{equation}
    \Delta L(P_{50})=755.8\,\mathrm{ms}-571.1\,\mathrm{ms}
                   =184.7\,\mathrm{ms}.
\end{equation}

\noindent The principal quantitative finding is therefore
\emph{a successful-flow median displacement of
$\Delta L(P_{50})=184.7$~ms}. This value places the centre of the observed
mediated distribution to the right of the direct comparison distribution. It
does not imply that every request incurs the same displacement.

\noindent The interpretation remains constrained by the implemented flow definitions.
Chapter~\ref{chap:selected-approach} establishes that the direct patient paths
contain two steps, whereas the mediated patient workflow contains five steps
and different contracts. The subtraction consequently describes the declared
campaign populations; it does not constitute a single-variable estimate of a
semantically equivalent operation or an operational budget attributable only
to gateway mediation.

\noindent Network transit, connection and security handling, policy enforcement,
transformation, routing, queueing, downstream execution, retries, and the
different workflow compositions can all contribute to the aggregate
difference. The result remains informative as a client-boundary path
displacement while respecting the surrogate-observability invariant established
in Chapter~\ref{chap:stateofart}.

\noindent The positive median displacement also does not demonstrate execution
instability or stability on its own. Those interpretations require joint
inspection of centre, tail, success rate, and delivery drift within a specified
flow mix and load regime. The thesis therefore retains the measured
distributional difference while separating it from mediation-only and
component-level causal attribution.

\subsection{Asynchronous State Truncation and Fail-Fast Efficiency}
\label{subsec:conclusions-state-truncation}

\noindent The framework preserves a strict execution invariant: steps within one flow
occurrence execute sequentially, while independent occurrences execute
concurrently. Let an occurrence traverse the ordered states
$X_0\rightarrow X_1\rightarrow\cdots\rightarrow X_m$, and let the random
index $\tau$ identify the first unsuccessful transition. The execution engine
implements the stochastic truncation operator
\begin{equation}
    \mathcal{T}_{\tau}(X_0,\ldots,X_m)
      =(X_0,\ldots,X_{\tau}),
\end{equation}
so the realised workflow contains only the valid prefix.

\noindent If $D_j$ denotes the service demand associated with downstream state $j$,
the fail-fast rule suppresses
\begin{equation}
    W_{\mathrm{avoided}}(\tau)
       =\sum_{j=\tau+1}^{m}D_j.
\end{equation}
The operator is stochastic because the position and cause of the first failure
vary across occurrences. Once $\tau$ occurs, the stopping rule remains
deterministic and the suffix receives no dispatch.

\noindent The outcome-stratified path comparison supplies the empirical counterpart to
this model. Successful mediated flows exhibit a median path displacement of
$+184.7$~ms, whereas failed mediated flows exhibit
\emph{a median differential of $+12.88$~ms} relative to failed direct flows.
The mean and ninety-fifth-percentile failed-flow differentials, respectively
$+62.18$~ms and $+284.30$~ms, preserve the non-negligible tail of the failure
population.

\noindent Figure~\ref{fig:stochastic-timed-automaton} condenses this
outcome-conditioned contrast into the representative direct patient-query
state progression.

\begin{figure}[H]
    \centering
    \begin{tikzpicture}[
        scale=0.94,
        transform shape,
        >=Stealth,
        node distance=22mm and 26mm,
        every node/.style={align=center},
        workflowstate/.style={circle, draw=veryPeri,
            fill=veryPeriLight, thick, minimum size=19mm,
            text width=17mm, font=\footnotesize},
        terminalstate/.style={circle, draw=red!80,
            fill=red!10, thick, minimum size=19mm,
            text width=17mm, font=\footnotesize},
        transition/.style={->, thick},
        failure/.style={->, thick, red!80}
    ]
        \node[workflowstate] (s0) {$S_0$\\\textbf{Idle}};
        \node[workflowstate, right=of s0] (s1)
            {$S_1$\\\textbf{Assertion}\\\textbf{acquired}};
        \node[workflowstate, right=of s1] (s2)
            {$S_2$\\\textbf{Workflow}\\\textbf{complete}};
        \node[terminalstate, below=15mm of s1] (sf)
            {$S_{\mathit{fail}}$\\\textbf{Early}\\\textbf{exit}};

        \draw[transition] ($(s0.west)+(-6mm,0)$) -- (s0.west);
        \path[transition]
            (s0) edge node[above, font=\scriptsize]
                {[RVE-1.b]} (s1)
            (s1) edge node[above, font=\scriptsize]
                {[RVE-54]}
                node[below, font=\scriptsize]
                {$\Delta L(P_{50})=184.7\,\mathrm{ms}$} (s2);
        \path[failure]
            (s0) edge[bend right=15]
                node[left, font=\scriptsize]
                {Unsuccessful\\transition $\tau$} (sf)
            (s1) edge node[right, font=\scriptsize]
                {$\mathcal{T}_{\tau}$\\$W_{\mathrm{avoided}}$} (sf);

        \node[draw=gray, dashed, rounded corners,
            fit=(s1)(s2), inner sep=4mm,
            label={[gray, font=\tiny]above:Protected downstream path}] {};
        \node[red!80, font=\scriptsize, below=3mm of sf]
            {Failed-population median differential: $+12.88\,\mathrm{ms}$};
    \end{tikzpicture}
    \caption[Stochastic timed automaton of workflow execution]{Empirically
    annotated workflow automaton representing transaction-state progression,
    stochastic truncation $\mathcal{T}_{\tau}$, and early execution
    termination. Source: Author's elaboration.}
    \label{fig:stochastic-timed-automaton}
\end{figure}

\noindent The contrast between $+12.88$~ms and $+184.7$~ms is consistent with a
shorter realised prefix for at least half of the interrupted occurrences. It
does not state that failure occurs within $12.88$~ms: the measured quantity is
the median mediated-minus-direct differential within the failed-flow
population.

\noindent The normal branch in Figure~\ref{fig:stochastic-timed-automaton} advances
from $S_0$ to $S_1$ when [RVE-1.b] produces the Security Assertion Markup
Language assertion required by the protected call. A successful [RVE-54]
transition then reaches $S_2$. Its $184.7$~ms annotation reports the median
displacement of the successful path populations; it is neither a transition-time
bound nor a probability attached to the automaton.

\noindent An unsuccessful prerequisite instead reaches $S_{\mathit{fail}}$ and
suppresses the remaining dependency suffix. For flow families that contain
document transactions, this branch prevents semantically invalid Integrating
the Healthcare Enterprise IT Infrastructure transaction 18 or transaction 43
requests from entering the downstream path. When the security prerequisite
already fails, those unexecuted requests also avoid the run-specific
$45$-second step-timeout regime. The figure assigns no transition probability
because the campaign reports outcome-conditioned populations rather than a
fitted probability for one common automaton.

\noindent Stochastic truncation connects the workflow state model to the load-state
model through avoided demand without conflating their state spaces. The former
describes dependency progression inside one healthcare occurrence; the latter,
adapted from Akram et al., relates traffic intensity, latency, and the validity
of distributed state under load\footnote{%
\phantomsection\label{fn:akram-load-model}%
\cite{akram-numerical-latency}.}.

\noindent At the gateway boundary, suppression of an invalid suffix removes requests
that would otherwise compete for admission, connection handling, parsing,
validation, transformation, and routing. At the downstream boundary, services
do not receive requests whose inputs depend on a failed state. This reduction
in submitted work can weaken one path to queueing backpressure, but the
campaign does not isolate its marginal effect on congestion. The mechanism
therefore establishes fail-fast termination as an execution and correctness
property; it neither reveals gateway capacity nor eliminates transfer in vain
for requests that already enter a long queue or approach a timeout.

\subsection[Queueing Saturation and Load-Sharing Window Closure]{Queueing Saturation and Load-Sharing Window}
\label{subsec:conclusions-queueing-saturation}

\noindent Chapter~\ref{chap:introduction} asks whether the experimental method exposes
behaviour at extreme offered load rather than only nominal latency. The
calibration evidence in Chapter~\ref{chap:selected-approach} locates the
relevant transition at the client-side admission mechanism, when
\emph{active flows reach the configured capacity $C_{\max}=30$ at approximately
$60$~s}. The scheduling-drift envelope then passes approximately
$60{,}000$~ms and approaches $67{,}000$~ms before the arrival stream drains.
Flow occurrences continue to arrive, but the executor cannot admit them at
their planned instants; delivery drift therefore accumulates outside the
semaphore.

\noindent The plateau identifies the generator's configured concurrency boundary, not
a regional service capacity of $30$. Once the semaphore remains occupied, the
planned arrival process $\lambda(t)$ no longer equals the workload delivered to
the remote path. Request latency excludes the interval spent waiting for local
admission, whereas scheduling drift records the difference between planned and
actual dispatch. Reading the two measures together therefore separates offered
work from admitted work and prevents late dispatches from being reclassified as
normal arrivals.

\begin{figure}[htpb]
    \centering
    \resizebox{0.70\textwidth}{!}{%
    \begin{tikzpicture}[x=0.9cm, y=0.9cm, font=\footnotesize, >=Stealth]
        \node[anchor=south west, yshift= 5mm, inner sep= 0pt, font=\footnotesize\bfseries] at (0.4,4.8)
            {Planned stochastic arrivals $t_{\mathit{plan}}$};
        \draw[->, thick] (0.4,4.3) -- (10.2,4.3);
        \foreach \x/\lab in {1.3/1,3.0/2,4.7/3,6.4/4,8.1/5} {
            \fill[veryPeri] (\x,4.3) circle (2pt);
            \node[veryPeri, above=0.1mm] at (\x,4.3) {$\tau_{\lab}$};
        }

        \draw[->, thick] (0.4,2.4) -- (10.2,2.4);
        \node[anchor=north west, font=\footnotesize\bfseries] at (0.4,2.25)
            {Observed saturated asynchronous timeline $t_{\mathit{actual}}$};
        \fill[teal] (1.4,2.4) circle (2pt);
        \fill[teal] (3.8,2.4) circle (2pt);
        \fill[teal] (6.8,2.4) circle (2pt);
        \draw[->, thick, teal] (1.3,4.1) -- (1.4,2.6)
            node[midway, right, font=\tiny] {$\delta_1$};
        \draw[->, thick, teal] (3.0,4.1) -- (3.8,2.6)
            node[midway, right, font=\tiny] {$\delta_2$};
        \draw[->, thick, teal] (4.7,4.1) -- (6.8,2.6)
            node[midway, right, align=left, font=\tiny]
            {observed envelope\\$\approx60{,}000$--$67{,}000\,\mathrm{ms}$};

        \node[anchor=west, font=\footnotesize\bfseries] at (0.4,1.1)
            {Proposed bounded-release timeline $t_{\mathit{rt}}$};
        \draw[->, thick] (0.4,0.6) -- (10.2,0.6);
        \foreach \x in {1.34,3.04,4.74,6.44,8.14} {
            \fill[red!80] (\x,0.6) circle (2pt);
        }
        \node[red!80, anchor=north] at (5.0,0.45)
            {Proposed obligation: $J_i^{\mathit{dispatch}}\leq\epsilon_i$};
    \end{tikzpicture}
    }
    \caption[Timing chronogram of queueing backpressure]{Comparative timing
    chronogram separating the planned arrival timeline, observed delivery drift, and the bounded-release real-time dispatcher. Source: Author's
    elaboration.}
    \label{fig:scheduling-drift-chronogram}
\end{figure}

\noindent Figure~\ref{fig:scheduling-drift-chronogram} represents this distinction as a
chronogram. The upper timeline contains the planned occurrences generated from
$\lambda(t)$; the middle timeline shows their progressively delayed admission
under cooperative event-loop saturation. The lower timeline states the proposed future obligation for a regulated real-time
dispatcher: each release remains within a declared bound $\epsilon_i$.

\noindent The extreme-stress scenario adds the remote-path and completion evidence needed
for stochastic interpretation. Relative to the paired constant baseline, pooled
successful Flow Completion Time rises by $13.45\%$ at the fiftieth percentile,
$85.44\%$ at the ninety-fifth percentile, and $17.70\%$ at the ninety-ninth
percentile, while flow success decreases by $3.63$ percentage points. This
non-uniform movement produces a deformed, heavy-tailed distribution rather than
a uniform location shift.

\noindent The request population also develops a bimodal latency distribution with a
secondary concentration near the run-specific $45$-second step timeout.
Successful Integrating the Healthcare Enterprise IT Infrastructure transaction
18 flows reach $37{,}869.5$~ms at the ninety-fifth percentile and
$43{,}972.1$~ms at the ninety-ninth percentile. Their upper tail therefore lies
in the same tens-of-seconds regime, although the flow and request populations
remain distinct.

\noindent Akram et al. define the Load-Sharing Window (LSW) as the stochastic interval
within which state remains useful for distributed load transfer. As traffic
intensity and communication delay rise, this window narrows until a transfer
arrives too late to produce useful work, a condition termed Transfer in Vain
(TIV)\footref{fn:akram-load-model}.

\noindent The experiment does not observe a regional load-transfer protocol and does
not claim literal identity with that topology. It observes the homologous
application-level signature: admission stays saturated, delivery drift expands,
the latency tail grows, and the proportion of useful terminal completions falls.
These converging signals are consistent with a closing load-sharing window at
the client boundary.

\noindent The timeout-bound mode marks the practical transfer-in-vain risk boundary for
this workload. An occurrence that retains a local task, connection, buffer, or
semaphore position while its transaction approaches the configured $45$-second
timeout consumes transport and execution resources without guaranteeing useful
completion before the configured limit.

\noindent Transaction 18 tail values close to this limit show how client resources
remain committed to work that is at increasing risk of termination. The
evidence supports this risk interpretation, rather than a claim that every
timed-out request is provably futile at dispatch time. Taken together,
semaphore occupancy, delivery drift, outcome-conditioned latency, and useful
completion make the closing window empirically testable at the client boundary.

\section{Limitations of the Empirical Evidence}
\label{sec:conclusions-limitations}
\subsection{Overview}
\noindent The conclusions remain conditional on the observation perimeter, workload
construction, and execution platform of the campaign. These boundaries do not
invalidate the measurements; they define the population and causal claims to
which those measurements apply. Four limitations are particularly important
for interpreting the path-level overhead, the saturation signature, and the
external validity of the method.

\paragraph{Black-Box Observability Boundary}
\noindent The regional integration path remains a black box to the experiment. Internal
logs from the Application Profile Management System (APMS), gateway components,
and downstream services are inaccessible within the authorised security and
policy perimeter. The Request Dataset Generator (RDG) therefore measures the
interval between client dispatch and client receipt and reconstructs state from
client-visible outcomes.

\noindent The method measures a declared path differential, but it cannot partition
that differential among network transport, connection establishment, security
processing, queue residence, transformation, routing, gateway computation, and
downstream execution. Consequently, \emph{$\Delta L(P_{50})=184.7$~ms remains
an aggregate surrogate displacement}, not a direct measurement of gateway
service time.

\noindent A causal decomposition requires aligned ingress, component-level,
downstream, and egress timestamps under a common correlation identifier and a
quantified clock-error bound. Without that instrumentation, any allocation of
the displacement to one internal phase remains unsupported.

\paragraph{Workload and Semantic Scope}
\noindent The workload reflects the Veneto Sistema Informativo Territoriale (SI.Ter)
integration perimeter represented in the available specifications and current
implementation. It exercises Anagrafe Zero patient interactions, Simple Object
Access Protocol (SOAP) document transactions including ITI-18 and ITI-43, and
Scryba Sign operations governed by the available developer guide
\footcites{arsenal-anagrafe-zero-v26,affinity-domain-italia-v263}
{scryba-sign-developer-guide-v16}.

\noindent These profiles make the experiment domain-specific at the level of the
represented semantics: security prerequisites, patient identities, document
references, and signature outputs remain ordered rather than collapsing into
independent endpoint calls. They do not establish operational
representativeness or define a universal healthcare workload.

\noindent The measured stochastic distributions cannot be transferred without
validation to purely Representational State Transfer (REST) systems, GraphQL
interfaces, message-broker topologies, or fully decoupled event-driven
microservices. Those architectures change connection reuse, synchronisation,
backpressure, state ownership, and completion semantics.

\paragraph{Local Client Ingestion Bottlenecks}
\noindent The extreme campaign also profiles the limits of a single load-generator node.
Once active flows occupy the semaphore, queued occurrences accumulate delivery
drift and the realised input to the remote system diverges from the configured
arrival process. The observed plateau and drift therefore demonstrate
\emph{local admission saturation}, but they do not isolate one machine-level
cause.

\noindent Event-loop scheduling, operating-system context switching, socket
availability, memory pressure, logging, and flow-level semaphore waiting can all
contribute. This confounding effect limits the offered rate $\lambda$ that the
experiment can attribute exclusively to the system under test.

\noindent Above the local knee, an increase in scheduled traffic can increase
client-side queueing without increasing remotely delivered traffic by the same
amount. The extreme run consequently characterises the coupled generator--path
system after local saturation rather than an isolated provider capacity.

\paragraph{Replication and Population Balance}
\noindent A cross-cutting sample limitation reinforces these three boundaries. The final
evidence set contains $74$ complete live runs from a protocol that specifies
$300$, and the number of complete repetitions differs across scenarios. The
large number of request and flow occurrences improves descriptive precision for
the observed populations, but it does not replace balanced independent
replication.

\noindent The campaign therefore supports empirical characterisation and model
validation within its realised sample. It does not establish a universal
capacity threshold, a service-level guarantee, or production-wide traffic
representativeness.

\section{Future Directions}
\label{sec:conclusions-future-work}

\subsection{Engineering Scalability and Platform Evolutions}
\label{subsec:future-engineering-scalability}

\noindent The engineering roadmap separates generator capacity from gateway
capacity by scaling, controlling, and grounding the experimental instrument.
The three evolutions below address, respectively, distributed workload
generation, closed-loop load control, and trace-based workload construction.
Each retains the thesis's evidence discipline by making calibration conditions,
observation boundaries, and failure states explicit.

\paragraph{Distributed Workload Generation}
\noindent The first evolution replaces the single-node Request Dataset Generator
(RDG) with a Kubernetes-orchestrated cluster whose workers execute disjoint
partitions of one global occurrence space. A control plane assigns deterministic
seed intervals, immutable configuration digests, globally unique occurrence
identifiers, and a common experiment epoch. Each worker maintains a local
asynchronous pool and publishes occurrence-scoped events to a partitioned
collector. If worker $k$ delivers intensity $\lambda_k(t)$, the intended offered
load is
\[
    \lambda(t)=\sum_{k=1}^{n}\lambda_k(t).
\]
This identity is accepted empirically only when per-worker scheduling drift
remains below a calibration bound and remotely observed arrivals grow
proportionally with $n$.

\noindent A two-phase start barrier separates image and connection warm-up from
timed execution. Monotonic clocks measure local durations, while a
clock-synchronisation service supplies an uncertainty envelope for cross-node
ordering. The manifest records worker placement, clock error, image version,
seed range, and partial-failure state in addition to the present run metadata.
The design removes the single event loop as the immediate offered-load
bottleneck, supports geographically distinct entry points, and preserves the
runtime-level distinction between concurrency, distribution, and cloud
scalability\footcites{unipd-runtimes-concurrency-syllabus}
{apache-jmeter-remote-testing}.

\paragraph{Adaptive Closed-Loop Workload Control}
\noindent The second evolution turns the open-loop scheduler into a constrained
feedback experiment. At control instant $k$, the controller observes
\[
 y_k=\left(P_{95}^{\mathrm{drift}},P_{95}^{\mathrm{latency}},
 e_k,X_k,C_k\right),
\]
where $e_k$ is the error rate, $X_k$ is useful-flow throughput, and $C_k$ is
admitted concurrency. It manipulates either the next offered rate or the
admission bound. A Proportional--Integral--Derivative (PID) law can compute
\begin{equation}
    u_k=K_p\varepsilon_k
        +K_i\sum_{h=0}^{k}\varepsilon_h\Delta t
        +K_d\frac{\varepsilon_k-\varepsilon_{k-1}}{\Delta t},
\end{equation}
with saturation, anti-windup, cool-down intervals, and explicit request and
error budgets.

\noindent Change-point logic marks a candidate stochastic knee when
$\partial X/\partial\lambda$ approaches zero while the derivatives of tail
latency, delivery drift, or failures remain positive. Repeated approaches from
below, under independent seeds and a fixed semantic mix, produce a confidence
region for that knee rather than an unsupported exact gateway threshold. A
Reinforcement Learning agent can replace the fixed controller when the response
surface is history-dependent, but it trains on mock or isolated traces and
operates live only behind a safety shield that constrains its action space. The
controller also checks per-worker headroom, because a feedback law that observes
only end-to-end latency can misclassify generator saturation as remote
saturation.

\paragraph{Anonymised Real-World Trace Replay}
\noindent The third evolution constructs an authorised trace-ingestion and replay
pipeline for regional clinical traffic. Parsers map heterogeneous operational
records to a canonical schema; privacy transformations remove unnecessary
payloads, tokenise identifiers consistently within one authorised trace,
generalise rare attributes, and shift absolute timestamps while preserving
inter-arrival intervals. A dependency-reconstruction stage then recovers causal
order, retries, flow mix, payload-size classes, and prerequisite relationships
before compiling the trace into RDG occurrences and a non-stationary arrival
timeline $\lambda(t)$.

\noindent Replay fidelity is assessed through approved aggregate properties---
including autocorrelation, burst duration, concurrency, and flow proportions---
and not through access to identifying clinical content. A segment
whose initial state or dependencies remain incomplete is rejected rather than
converted into an independent request stream, because accurate scaling preserves
both temporal and causal structure\footcite{zhu-tbbt-trace-replay}. The pipeline
keeps the re-identification secret outside the experimental environment, records
an auditable transformation manifest, and applies data minimisation and
disclosure-risk review before any replay artefact enters the cluster
\footcite{eu-gdpr-2016-679}.

\subsection{Theoretical Computer Science Perspectives}
\label{subsec:future-theoretical-perspectives}

\noindent The theoretical roadmap converts observed regularities into proposed proof obligations. These directions extend the
empirical contribution; they do not report proofs completed by this thesis.
Each obligation therefore remains conditional on an explicit abstraction, an
identified population, and a declared observation boundary.

\paragraph{Predictable Dispatch and Release Jitter}
\noindent The first proposed extension treats the divergence represented in
Figure~\ref{fig:scheduling-drift-chronogram} as a local schedulability problem.
The non-stationary Poisson generator admits arbitrarily short inter-arrival
distances, so a regulator first maps each flow class to a sporadic task
$\tau_i=(C_i,T_i,D_i,J_i)$ with worst-case local demand $C_i$, minimum
inter-arrival time $T_i$, deadline $D_i$, and release jitter $J_i$. Under the
classical independent, implicit-deadline assumptions, the sufficient
utilisation bounds for Rate Monotonic and Earliest Deadline First scheduling on
one processor are, respectively,
\[
 \sum_i\frac{C_i}{T_i}\leq n\left(2^{1/n}-1\right),
 \qquad
 \sum_i\frac{C_i}{T_i}\leq1.
\]
Shared resources require the response-time recurrence
\begin{equation}
 R_i=C_i+B_i+\sum_{j\in hp(i)}
       \left\lceil\frac{R_i+J_j}{T_j}\right\rceil C_j,
 \qquad R_i+J_i\leq D_i,
\end{equation}
where $B_i$ bounds local blocking and $hp(i)$ contains the higher-priority
tasks.

\noindent The analytical target for the lower chronogram is
$J_i^{\mathit{dispatch}}=\max(0,s_i-r_i)\leq\epsilon_i$, with release time
$r_i$, actual start $s_i$, and a declared engineering bound $\epsilon_i$; it is
not a measured microsecond guarantee. A qualified hard real-time platform can host the timing-critical dispatcher, while
the \texttt{PREEMPT\_RT} real-time-preemption configuration provides an
intermediate Linux option. The Priority Ceiling Protocol can bound inversion in
local credential stores, connection pools, and log buffers used by mutual
Transport Layer Security and Security Assertion Markup Language processing.
These mechanisms bound local scheduling effects, not the remote network
itself\footcites{unipd-real-time-kernels-syllabus}
{liu-real-time-systems}{burns-wellings-analysable-real-time}.

\paragraph{Static Verification of Concurrency Bounds}
\noindent The second proposed extension applies abstract interpretation to the
asynchronous executor's lowered control-flow graph. Let $\Sigma$ be the concrete
state space and let $F(X)=\Sigma_0\cup\operatorname{post}(X)$ be the monotone
collecting transformer. Under the required $\omega$-continuity assumption, its
concrete collecting semantics is
\begin{equation}
 \mathcal{S}_{\mathit{conc}}=\operatorname{lfp}(F)
   =\bigcup_{n\in\mathbb{N}}F^n(\varnothing).
\end{equation}
A Galois connection
\begin{equation}
 \left\langle\mathcal{P}(\Sigma),\subseteq\right\rangle
 \underset{\gamma}{\overset{\alpha}{\rightleftarrows}}
 \left\langle D^{\sharp},\sqsubseteq\right\rangle
\end{equation}
maps concrete states to an interval domain
$D_{\mathit{Int}}^{\sharp}$ and a relational octagon domain
$D_{\mathit{Oct}}^{\sharp}$.

\noindent A sound transformer $F^{\sharp}$ computes a
post-fixpoint through widening,
\begin{equation}
 X_{k+1}^{\sharp}=X_k^{\sharp}\nabla F^{\sharp}(X_k^{\sharp}),
\end{equation}
and aims to infer, for active flows $a$, held permits $h$, and available permits
$p$, the inductive invariant
\begin{equation}
 0\leq a\leq h\leq30,\qquad 0\leq p\leq30,\qquad h+p=30.
\end{equation}
If the analyser establishes this invariant for normal completion, exceptions,
and cancellation, soundness yields $a\leq30$ for every arrival trace.

\noindent This proof remains future work. It also does not bound waiting tasks:
the present scheduler can create an unbounded waiting population, so proving
$0\leq q\leq Q$ additionally requires a bounded producer queue and an explicit
block, shed, or defer policy
\footcites{unipd-software-verification-syllabus}{mine-numeric-invariants}.

\endgroup
\newpage
